\section{Adaptive planar enumeration of complete standard rays}
\label{sec:planar-sectors-native}

Report 80 extends the minimum-envelope method from a projective interval
to a projective polygon. Automatic STANDARD enumeration now selects it
when matching nullity is three. The existing nullity-one/two envelope
paths remain selected, and higher-dimensional sectors retain the earlier
support/arrangement choice. Explicit planar enumeration also accepts
lower nullity. Selection is based on an exact geometric parameter rather
than elapsed-time guesses. Complete sector discovery still uses its
existing independently replayable disc/exhaustion consumer.

\subsection{The common minimum fan describes actual standard faces}

Let $k$ be the number of allowed labelled quadrilateral variables and $d$
the dimension of their matching kernel. Intersect its nonnegative cone
with $\sum q_i=1$. This compact section $P$ has effective dimension at most
$d-1$; raw nullity three can still give an empty set, point or segment.
For each global vertex the triangle potentials are linear forms, and the
canonical lift subtracts their minimum. Choosing an attaining corner
therefore defines a convex minimum cell in $P$, by inequalities comparing
its form to all the others. Identical projected forms are collapsed.

The common refinement of these minimum fans and the coordinate faces of
$P$ describes the canonical standard coordinate faces. A point is an
extreme standard ray precisely at a vertex of this refinement. This
follows directly from the matching equations and zero-coordinate faces:
within a fixed minimum cell the canonical lift is linear; conversely,
the zero triangle and quadrilateral coordinates defining a standard face
recover the corresponding minimum and coordinate-face constraints.
Thus Q-cone extremes alone need not give all standard extremes.

In a two-dimensional section, full-dimensional minimum cells and their
faces suffice. A label active only on a boundary face adds no geometry:
the difference between it and the adjacent minimum is a nonnegative linear
function, whose zero set on a face cannot be an isolated interior point.
The implementation clips these convex cells, extracts their internal
segments, and overlays segments from different global-vertex groups.
The rays are the original polygon vertices, individual fan vertices and
intersections between group segments. No standard-rank filter is needed
to remove spurious arrangement crossings.

\begin{proposition}[Output and construction bounds]
Let $b$ be the number of vertices of $P$, and $m_v$ the number of distinct
full-dimensional active minimum cells at global vertex $v$. Then
\[
 R\le b+2\sum_v(m_v-1)
           +9\sum_{v<w}(m_v-1)(m_w-1)=O(k^2).
\]
With one nontrivial minimum group the bound is linear. After kernel
construction, a deterministic realization has elementary scalar-operation cost
$O(k^3+tk^2)$; rational coefficient-bit costs must also be included.
\end{proposition}
\begin{proof}
A convex subdivision into $m$ minimum cells has at most $2(m-1)$ new
vertices and at most $3(m-1)$ internal segments. Euler's planar formula
and degree at least three at genuine new vertices give the first bound;
the dual planar adjacency bound gives the second, since two convex cells
share at most one segment. Summing individual fan vertices and the at most
one crossing for each pair of segments from different groups gives the
displayed output bound. The source equations provide only $O(k)$ relevant
distinct potential forms after the equality contraction.
Clipping every cell against its group's dominance inequalities costs
$O(k^3)$ in this elementary construction; overlay costs $O(k^2)$ and lifting
at most $O(k^2)$ rays into $7t$ coordinates costs $O(tk^2)$.
The exact rational charts and intersections have polynomially bounded bit
length in source coefficient bits. A unit-cost rational count alone would
not establish the bit-complexity claim.
\end{proof}

\subsection{Independent area and length coverage replay}

\begin{sloppypar}
The producer is \code{sector\_planar.py}; the independent checker is
\code{sector\_planar\_verify.py}, using the source-bound
\code{normal-sector-planar-v1} format. Its source model reconstructs the
full uncollapsed corner matching graph and checks the chart and bounds.
It does not import the producer's polygon clipper. Each submitted convex
cell has source-checked vertices and dominance inequalities. Linear
dominance at all vertices proves dominance throughout the cell.
\end{sloppypar}

Within a group, cells with different projected active forms have disjoint
interiors: simultaneous equality on an open set would make the forms
identical, contrary to canonical collapse. The checker verifies that their
exact areas sum to the area of $P$. A finite union of closed cells cannot
leave an open gap of positive area, and closure then covers its boundary.
Hence the cells cover $P$. For effective dimension one the corresponding
argument uses exact lengths; points and empty sections have separate
controls. The checker independently reconstructs the overlay vertices and
compares every normalized primitive lifted ray with the submitted list.
Missing rays, malformed fractions, incorrect charts, hidden lower forms,
source substitutions and cancellation are rejected.

This proves coverage of one supplied compatible sector. A negative disc
result also checks every ray's disc predicate independently. It does not
prove that the selected sector family contains a disc for every unknot,
and its status must not be promoted to a knottedness verdict.

\subsection{Reusable sparse source preparation, kept explicit}

The additive \code{sector\_sparse.py} module validates and prepares immutable
source geometry once. Building another sector reuses that prepared source,
whose identity is checked before certificate emission. Union by size and
path halving contract zero-type corner equalities. Only $4k$ active labelled
edges carry dense length-$k$ potential tuples; zero records share constants.
Fundamental-cycle coefficients lie in $[-4,4]$, permitting deterministic
radix ordering in $O(k^2)$. Pre-elimination work is
$O(t\log t+k^2)$, with subsequent rational elimination accounted separately.
This reusable interface is valuable for repeated sector queries.
The delivered one-shot measurements are slower than the current dense
constructor, so native automatic one-shot enumeration retains that
constructor. The planar enumeration benchmark below charges both arms
for the same native full kernel preparation.

\subsection{Seven rays and a uniform integral disc classification}

The report gives a geometric family with $n$ Fibonacci base tetrahedra
and two boundary caps. Use base type two and cap type one, and put
$A=F_{n+1}$, $B=F_{n+2}$, $F_0=0$, $F_1=F_2=1$.
Every standard matching solution has unique parameters $x,y,z,b,u,v$:
\begin{align*}
 i<n:\quad &(F_{n-i+1}x+b,F_{n-i+1}x+b,b,b;0,0,F_{n-i}x),\\
 i=n:\quad &(Ax+b-y,Bx+b,u,b;0,y,0),\\
 i=n+1:\quad &(Bx+b,Ax+b-z,b,v;0,z,0).
\end{align*}
Nonnegativity is exactly $x,b,u,v\ge0$ and $0\le y,z\le Ax+b$.
The matching kernel has dimension three. Canonical lifting sets
$h=\max(0,y-Ax,z-Ax)$, $b=h$, $u=v=0$; general solutions add the three
boundary vertex links. Substitution into the base equations gives
$W_i=W_{i+1}+W_{i+2}$ and hence the Fibonacci rows, while the two attaching
face equations determine all cap triangles except $u,v$.

\begin{proposition}[All-size double-cap classification]
For every $n\ge1$, the canonical sector has precisely the seven rays
\[
 (1,0,0),\ (1,A,0),\ (1,0,A),\ (1,A,A),\
 (0,1,0),\ (0,0,1),\ (0,1,1).
\]
The first four are essential meridian discs; the last three are inessential
discs. Every integral canonical vector consists of $x$ meridian discs
and $h$ inessential discs. A full integral vector consists of $x$ meridian
discs and $b+u+v$ inessential discs.
\end{proposition}
\begin{proof}
The three active forms are $0,Ax-y,Ax-z$. Their minimum fan in the
nonnegative projective triangle has three cells and exactly the seven
displayed vertices. Coordinate-face extremality proves these are all rays.
For components, the base restriction is $x$ parallel meridian discs and
$b$ vertex-link discs, identified by uniqueness of normal coordinates.
With $u=v=0$, every cap piece meets its attaching face in one arc; the
opposite-vertex triangle is absent. Attaching a disc leaf along one interval
preserves a disc and cannot merge components. Apply this to the two caps
successively. The arcs use distinct global edges even when attaching-face
vertices are identified in the quotient.

Essentiality is preserved by replacing each old arc with the corresponding
cap-boundary arc. Two disjoint systems in an abstract disc with identical
endpoint pairing are related by a PL homeomorphism fixed on the boundary,
constructed by cutting outermost arcs. It respects boundary identifications
and descends to the quotient boundary torus. Thus the original meridians
remain essential and links remain inessential. The free triangles $u,v$
are the new vertex links. Substituting $b=h$ gives the canonical census.
The preserved realization proof handles quotient faces with an interior
positive collar vanishing on the boundary, rather than assuming embedded
face closures.
\end{proof}

Only $(1,0,0)$ among the essential rays is a Q-cone extreme; three further
essential standard rays are created by the minimum fan. This uniform
theorem explains why the planar enumeration is needed. Primitive integral
coordinates alone still do not guarantee a connected surface.

\subsection{Native validation and complete measurements}
All 1,391 maintained tests pass in 250.052 seconds. The focused 72-test
run passes in 2.205 seconds. Integration tests disable pairwise hyperplane
generation and standard-rank filtering in the planar branch; positive
discovery, independent producer-disabled coverage, missing-ray mutations,
shared-source identity and actual-ray caps have explicit controls.

The audit selects 9,995 sectors across 48 frozen triangulations and checks
every selected one of matching nullity at most three: 9,807 queries,
including 463 of nullity three. This selection is exhaustive only on the
declared smaller sources; it is not all supports on every large source.
All 7,674 ray occurrences match the frozen complete standard lists, including
2,279 non-Q rays. The nullity-three subset contributes 2,821 rays and
1,475 non-Q rays, with 129 essential-disc occurrences, 84 non-Q.
Automatic dispatch agrees with direct planar output on every audited
query. On all 9,344 lower-nullity queries, the complete output order and
statistics exactly equal the frozen incumbent. Eight fresh Regina complete
standard enumerations match the frozen lists. The 188 selected sectors of
higher raw nullity remain explicitly outside the gate.

Four additional fresh geometric families at base sizes 1, 2, 4 and 8
verify 300 full-coordinate parameter cases, forty independently replayed
component censuses, 28 fresh Regina ray classifications and four fresh
full standard enumerations. Every family has seven rays, four essential
and three inessential. These finite controls support the separate
all-size written proof; they do not substitute for it.

The benchmark compares the frozen native sector implementation at
\code{6a34c5c0d}. Both arms pay full native dense kernel construction.
Five interleaved four-arm rounds complete 220 measured calls and 44 warmups
on an otherwise idle observed host, with A/A arms retained. Complete
enumeration includes every output coordinate, serialization and ray digest.
That iterator has no portable coverage proof; such replay is not claimed
inside enumeration time. Complete discovery includes its positive or
exhaustion proof, independent verification and full proof serialization.
Both arms return the same discovery status and their actual witnesses are
verified. A smaller witness's differing byte length is recorded.

\begin{samepage}
\noindent\textbf{Complete standard enumeration and output.}
\begin{center}
\small
\begin{tabular}{lrrrrr}
\toprule
Input & old ms & new ms & old/new & old A/A & new A/A\\
\midrule
Base 1, two caps & 4.332 & 3.178 & 1.363 & 0.997 & 1.004\\
Base 4, two caps & 61.947 & 6.291 & 9.850 & 0.890 & 1.004\\
Base 8, two caps & 676.910 & 11.143 & 60.054 & 0.999 & 1.011\\
Base 16, two caps & 11931.837 & 24.230 & 493.136 & 1.008 & 1.001\\
Base 4, one cap & 1.639 & 1.620 & 1.008 & 1.001 & 0.999\\
Base 8, one cap & 3.522 & 3.537 & 0.992 & 0.993 & 1.012\\
\bottomrule
\end{tabular}
\end{center}

\end{samepage}

At base size 16, complete enumeration drops from about 11.932 seconds to
0.024 seconds, a paired 493.14-fold gain, producing the same seven rays.
Nullity-two one-cap controls retain the old envelope path and stay near
parity. The size-four old A/A ratio is 0.890 in the first run. A separate
fifteen-round repeat completes sixty measured calls and four warmups;
old and new A/A ratios are 1.011 and 1.001, with paired gain 9.823.
Both records are retained, rather than discarding the noisier run.

\begin{samepage}
\noindent\textbf{Complete supplied-sector disc discovery and replay.}
\begin{center}
\small
\begin{tabular}{lrrrrr}
\toprule
Input & old ms & new ms & old/new & old A/A & new A/A\\
\midrule
Base 1, two caps & 3.119 & 7.089 & 0.449 & 0.954 & 1.035\\
Base 4, two caps & 11.535 & 12.757 & 0.910 & 0.999 & 0.994\\
Base 8, two caps & 38.402 & 20.807 & 1.854 & 0.989 & 0.997\\
Base 16, two caps & 191.528 & 41.348 & 4.683 & 1.055 & 1.008\\
Empty sector & 1.384 & 1.410 & 0.989 & 0.985 & 1.020\\
\bottomrule
\end{tabular}
\end{center}

\end{samepage}

At size 16, complete discovery drops from about 0.192 to 0.041 seconds,
a paired 4.68-fold gain. It improves 1.85-fold at size eight, but is
about 2.23 times slower at size one and 10 percent slower at size four.
The old iterator can find its first disc before paying full arrangement
work, whereas the planar iterator prepares its minimum diagram before
yielding. These regressions limit the discovery policy; enumeration gains
must not be substituted for them. Callers can still select the old
support/arrangement methods explicitly. Full recognition is not timed by
these supplied-sector rows. The empty-sector control is near parity.

Four new-only complete coverage calls at base sizes 1, 8, 16 and 32
include production, independent coverage replay and serialization. Times
are approximately 0.008, 0.024, 0.054 and 0.182 seconds; proof sizes are
995, 2,048, 3,442 and 6,869 bytes. These are capacity observations without
a baseline ratio. Publication review separately replays the thirteen
native audit proofs, four family proofs, four capacity proofs and seventeen
delivered proofs with the producer disabled.

All 587 captured native source pins stay fixed during audit and timing.
They refer to runtime release \code{cc44aae87}; a subsequent docstring-only
edit documents the new dispatch and is checked to preserve its executable
syntax tree. Historical edge-first evidence retains its own release pins.
Reproduction uses \code{normal\_orbit\_research.planar\_sectors} with
\code{audit --fresh-regina} or \code{benchmark --rounds 5}, from \code{fast/}.
Raw evidence, the noisy-row repeat and publication checks are retained
under \code{data/planar-sectors-*}.

